
%%%%%%%%%%%%%%%%%%%%%%%%%%%%%%%%%%%%%%%%
%           main body
%%%%%%%%%%%%%%%%%%%%%%%%%%%%%%%%%%%%%%%%

%-----------------------------------
%	TITLE PAGE
%-----------------------------------

\title[数学分析]{\texorpdfstring{\LARGE \bfseries 第十章\quad 函数项级数}{第十章 函数项级数}} % The short title appears at the bottom of every slide, the full title is only on the title page
\author[函数项级数] % Your institution as it will appear on the bottom of every slide, maybe shorthand to save space
{\Large \bfseries 函数项级数的一致收敛性}
% \author{Singular Value Decomposition} % Your name
% \date{\today} % Date, can be changed to a custom date
\date{}

%------------------------------------------------

\begin{document}

%------------------------------------------------

\begin{frame}

\titlepage % Print the title page as the first slide

\end{frame}

%------------------------------------------------

% \begin{frame}
% \frametitle{内容提要} % Table of contents slide, comment this block out to remove it
% \tableofcontents % Throughout your presentation, if you choose to use \section{} and \subsection{} commands, these will automatically be printed on this slide as an overview of your presentation
% \end{frame}

%-----------------------------------
%	PRESENTATION SLIDES
%-----------------------------------

%------------------------------------------------

\begin{frame}{极限运算能否任意交换次序?}

设每个 $u_n(x)$ 在 $[a, b]$ 上连续, 函数项级数 $\sum\limits_{n=1}^\infty u_n(x)$ 在 $[a, b]$ 上逐点收敛于 $S(x).$

\begin{question}
$S(x)$ 是否也连续? 即能否有
\[
\lim_{x \to x_0} \sum_{n=1}^\infty u_n(x) = \sum_{n=1}^\infty \lim_{x \to x_0} u_n(x)?
\]
\end{question}

\pause
\vspace{0.4em}

\alert{答案是否定的.} 以函数列 $S_n(x) = x^n$ 在 $[0, 1]$ 为例:
\[
S(x) = \lim_{n\to\infty} x^n = \begin{cases} 0, & x \in [0, 1), \\ 1, & x = 1. \end{cases}
\]
每个 $x^n$ 在 $[0, 1]$ 上连续, 但极限函数 $S(x)$ 在 $x = 1$ 处间断!

\pause
\vspace{0.4em}

\alert{原因在于:} 对不同的 $x,$ $x^n \to 0$ 的速度差别悬殊: 靠近 $x = 1$ 处, 收敛越来越慢, 导致``收敛不一致''.

\end{frame}

%------------------------------------------------

\section[一致收敛]{一致收敛的概念}

%------------------------------------------------

\begin{frame}{一致收敛的定义}

\begin{Def}[一致收敛]
称函数项级数 $\sum\limits_{n=1}^\infty u_n(x)$ 在 $D$ 上\textbf{一致收敛}到 $S(x),$ 是指:
\[
\forall ~ \varepsilon > 0, ~ \exists ~ N(\varepsilon) ~ (\text{只与 } \varepsilon \text{ 有关}), ~ \text{使得} ~ \forall ~ n \geqslant N, ~ \forall ~ x \in D: ~ \left| S_n(x) - S(x) \right| < \varepsilon.
\]
\end{Def}

\pause
\vspace{0.5em}

\begin{empheq}[box={\fancymath[colback=cyan!30, drop lifted shadow, sharp corners]}]{equation*}
\text{\bfseries\color{red!80} 逐点收敛: $N$ 与 $\varepsilon$ 和 $x$ 都有关; \quad 一致收敛: $N$ 只与 $\varepsilon$ 有关.}
\end{empheq}

\pause
\vspace{0.5em}

等价地: $\sum u_n(x)$ 在 $D$ 上一致收敛 $\Longleftrightarrow$ $\sup\limits_{x \in D} |S_n(x) - S(x)| \to 0.$

\pause

回到上例: $\sup\limits_{x \in [0,1)} |x^n - 0| = 1 \not\to 0,$ 故 $x^n$ 在 $[0, 1)$ 上\alert{不一致收敛}到 $0.$

\end{frame}

%------------------------------------------------

\section[判别法]{一致收敛的判别法}

%------------------------------------------------

\begin{frame}{Weierstra\ss{} 判别法}

\begin{thm}[Weierstra\ss{} 判别法]
设 $|u_n(x)| \leqslant a_n$ 对所有 $x \in D$ 和 $n \in \mathbb{N}$ 成立. 若数项级数 $\sum\limits_{n=1}^\infty a_n$ 收敛, 则函数项级数 $\sum\limits_{n=1}^\infty u_n(x)$ 在 $D$ 上一致 (且绝对) 收敛.
\end{thm}

\pause

\startproof[bg=yellow, fg=red]
对 $m > n > N(\varepsilon)$ (其中 $N$ 由正项级数 $\sum\limits_{n=1}^\infty a_n$ 的 Cauchy 收敛原理确定), 对所有 $x \in D:$
\[
\left| \sum_{k=n+1}^m u_k(x) \right| \leqslant \sum_{k=n+1}^m |u_k(x)| \leqslant \sum_{k=n+1}^m a_k < \varepsilon.
\]
由 Cauchy 收敛原理, $\sum\limits_{n=1}^\infty u_n(x)$ 在 $D$ 上一致收敛.
\qed

\pause
\vspace{0.5em}

\begin{eg}
设 $\sum\limits_{n=1}^\infty |a_n|$ 收敛, 则 $\sum\limits_{n=1}^\infty a_n \sin(nx)$ 在 $\mathbb{R}$ 上一致收敛. (取 $|a_n \sin(nx)| \leqslant |a_n|.$)
\end{eg}

\end{frame}

%------------------------------------------------

\begin{frame}{内闭一致收敛}

\begin{eg}
判断 $\sum_{n=0}^\infty x^n$ 在 $D = (-1, 1)$ 上的一致收敛性.
\end{eg}

\pause

在整个 $D = (-1, 1)$ 上, $\sup\limits_{x \in D} |x^n| = 1,$ 无法直接用 Weierstra\ss{} 判别法.

\pause
\vspace{0.5em}

实际上, 和函数 $S(x) = \dfrac{1}{1-x}$ 在 $x \to 1^-$ 时趋于无穷, 而各项 $|x^n| \leqslant 1$ 有界, 故 $\sum x^n$ 在 $(-1, 1)$ 上\alert{不一致收敛.}

\pause
\vspace{0.5em}

但对任意 $r \in (0, 1),$ 在 $D_r = [-r, r]$ 上有 $|x^n| \leqslant r^n,$ 而 $\sum r^n < \infty,$ 故 $\sum_{n=0}^\infty x^n$ 在 $(-1, 1)$ 上\alert{内闭一致收敛}: 在每个 $[-r, r] \subset (-1, 1)$ 上一致收敛.

\end{frame}

%------------------------------------------------

\section[分析性质]{和函数的分析性质}

%------------------------------------------------

\begin{frame}{和函数的连续性}

\begin{thm}[连续性定理]
设函数项级数 $\sum_{n=1}^\infty u_n(x)$ 在 $[a, b]$ 上一致收敛到 $S(x),$ 且每项 $u_n(x)$ 在 $[a, b]$ 上连续, 则 $S(x)$ 在 $[a, b]$ 上连续, 即极限运算与求和可以交换次序:
\[
\lim_{x \to x_0} \sum_{n=1}^\infty u_n(x) = \sum_{n=1}^\infty \lim_{x \to x_0} u_n(x), \quad \forall ~ x_0 \in [a, b].
\]
\end{thm}

\pause

\startproof[bg=yellow, fg=red]
固定 $x_0.$ 由三角不等式:
\[
|S(x) - S(x_0)| \leqslant \underbrace{|S(x) - S_n(x)|}_{\text{\Circled{1}}} + \underbrace{|S_n(x) - S_n(x_0)|}_{\text{\Circled{2}}} + \underbrace{|S_n(x_0) - S(x_0)|}_{\text{\Circled{3}}}.
\]
\pause
由\alert{一致收敛性,} 取 $N$ 使得对所有 $n \geqslant N$ 及所有 $x \in [a, b],$ $\text{\Circled{1}} < \varepsilon/3$ 且 $\text{\Circled{3}} < \varepsilon/3.$ 固定该 $n,$ 再由 $S_n$ 的\alert{连续性}取 $\delta > 0$ 使得 $|x - x_0| < \delta$ 时 $\text{\Circled{2}} < \varepsilon/3.$
\qed

\end{frame}

%------------------------------------------------

\begin{frame}{逐项积分定理与应用}

\begin{thm}[逐项积分定理]
设 $\sum_{n=1}^\infty u_n(x)$ 在 $[a, b]$ 上一致收敛到 $S(x),$ 且每项连续, 则积分与求和可以交换次序:
\[
\int_a^b S(x) ~ \mathrm{d}x = \sum_{n=1}^\infty \int_a^b u_n(x) ~ \mathrm{d}x.
\]
\end{thm}

\pause
\vspace{0.5em}

\begin{eg}
利用 $\dfrac{1}{1+t^2} = \sum_{n=0}^\infty (-1)^n t^{2n}$ 在任意 $|t| \leqslant r < 1$ 上内闭一致收敛, 对 $|x| < 1$ 逐项积分:
\[
\arctan x = \int_0^x \frac{\mathrm{d}t}{1+t^2} = \sum_{n=0}^\infty \frac{(-1)^n}{2n+1} x^{2n+1} = x - \frac{x^3}{3} + \frac{x^5}{5} - \cdots
\]
\end{eg}

\end{frame}

%------------------------------------------------

\section{总结}

%------------------------------------------------

\begin{frame}{总结}

\begin{itemize}
\item[{\larger[2] \color{red} \ding{43}}] \textbf{一致收敛的本质:} $N(\varepsilon)$ 与 $x$ 无关; 等价于 $\sup\limits_{x \in D} |S_n(x) - S(x)| \to 0.$

\pause
\vspace{0.5em}

\item[{\larger[2] \color{red} \ding{43}}] \textbf{Weierstra\ss{} 判别法:} $|u_n(x)| \leqslant a_n$ 且 $\sum\limits_{n=1}^\infty a_n < \infty$ $\Rightarrow$ $\sum\limits_{n=1}^\infty u_n(x)$ 在 $D$ 上一致收敛.

\pause
\vspace{0.5em}

\item[{\larger[2] \color{red} \ding{43}}] \textbf{分析性质 (一致收敛是充分条件):}
\begin{itemize}
\item 连续性: 各项连续 $+$ 一致收敛 $\Rightarrow$ 和函数连续;
\item 可积性 (逐项积分): 积分与求和可以交换次序;
\item 可微性 (逐项求导): 导函数的级数一致收敛时可逐项求导.
\end{itemize}
\end{itemize}

\pause
\vspace{0.5em}

\begin{attnblock}
一致收敛是上述性质的充分条件, 但并\alert{非}必要条件. 没有一致收敛性时, 上述性质也可能成立.
\end{attnblock}

\end{frame}

%------------------------------------------------

%------------------------------------------------
% % closing page
% \begin{frame}
% \HUGE{\centerline{\bfseries {\color{gray} The} {\color{zkblue} End}}}

% \pause

% \vspace{1.2em}

% \Large{\centerline{\bfseries {\color{gray}Thank} \quad {\color{zkblue} You}}}

% \end{frame}

%----------------------------------------------------------------------------------------

\end{document}
